\subsection{与滤过代数相伴的分次代数}

设 G 是一个带实滤过 \((G_{\alpha})_{\alpha \in \mathbb{R}}\) 的交换群。如同前面我们记

\[
\mathrm{G} _ {\alpha} ^ {+} = \bigcup_ {\beta > \alpha} \mathrm{G} _ {\beta};\tag{9}
\]

显然 \(\mathrm{G}_{\alpha}^{+}\) 是 \(\mathrm{G}_{\alpha}\) 的一个子群。我们记 \(\mathrm{gr}_{\alpha}(\mathrm{G}) = \mathrm{G}_{\alpha} / \mathrm{G}_{\alpha}^{+}\) ，并且

\[
\operatorname{gr} (\mathrm{G}) = \bigoplus_ {\alpha \in \mathbf {R}} \operatorname{gr} _ {\alpha} (\mathrm{G}).
\]

与滤过群 G 相伴的分次群是群 \(\operatorname{gr}(G)\) ，带其自然的 R 型分次。

注. 当滤过 \((\mathrm{G}_{\alpha})\) 是整滤过时， \(\mathrm{gr}_{\alpha}(\mathrm{G}) = \{0\}\) （对非整的 \(\alpha\) ），而 \(\mathrm{gr}_n(\mathrm{G}) = \mathrm{G}_n / \mathrm{G}_{n - 1}\) （对每个整数 \(n\) ）。因此相伴分次群的定义与《交换代数》第 III 章 § 2 no. 3 中的定义本质上是一致的。

设 \(\mathbf{A}\) 是一个代数（相应地，幺代数）， \((\mathbf{A}_{\alpha})_{\alpha \in \mathbb{R}}\) 是一个与代数（相应地，幺代数）结构相容的实滤过 (no. 1)。那么

\[
\mathrm{A} _ {\alpha}. \mathrm{A} _ {\beta} \subset \mathrm{A} _ {\alpha + \beta}, \quad \mathrm{A} _ {\alpha} ^ {+}. \mathrm{A} _ {\beta} + \mathrm{A} _ {\alpha}. \mathrm{A} _ {\beta} ^ {+} \subset \mathrm{A} _ {\alpha + \beta} ^ {+},
\]

，并且由 A 上乘法的限制所给出的 \(A_{\alpha} \times A_{\beta}\) 到 \(A_{\alpha+\beta}\) 中的双线性映射在取商后定义了一个双线性映射

\[
\operatorname{gr} _ {\alpha} (\mathrm{A}) \times \operatorname{gr} _ {\beta} (\mathrm{A}) \rightarrow \operatorname{gr} _ {\alpha + \beta} (\mathrm{A}).
\]

我们得到 \(\mathrm{gr}(\mathbf{A}) \times \mathrm{gr}(\mathbf{A})\) 到 \(\mathrm{gr}(\mathbf{A})\) 中的一个双线性映射，它使后者成为一个 R 型分次代数（相应地，幺分次代数）。若 A 是一个结合代数（相应地，交换代数，相应地，李代数），则 \(\mathrm{gr}(\mathbf{A})\) 亦然。

